\section*{Capítulo \ref{ch:g8:reproduction-and-environment} --- A reprodução e o ambiente}

\begin{solution}{exo:g8:reproduction-and-environment:1}
A regra: os filhotes precisam chegar quando o alimento e as \omterm{def:g6:exploring-our-environment:conditions}{condições}
os servirem --- para os filhotes de chapim, o pico de \omterm{ex:g2:animals-grow-up:butterfly}{lagartas} da
primavera nos carvalhos. A deixa que aciona a maquinaria dela: os dias
que se alongam.
\end{solution}

\begin{solution}{exo:g8:reproduction-and-environment:2}
\omterm{def:g2:born-grow-die:cycle}{Nascimentos} na primavera, quando o capim e o tempo ameno recebem os
filhotes. Os meses de gestação precisam ser contados para trás --- o
\omterm{prop:g4:animal-reproduction:where}{acasalamento} de outono é aritmética de primavera.
\end{solution}

\begin{solution}{exo:g8:reproduction-and-environment:3}
Porque cada \omterm{def:g2:animals-grow-up:born}{ovo}, cada \omterm{def:g2:from-seed-to-plant:seed}{semente} e cada gota de \omterm{def:g2:animals-grow-up:born}{leite} são \omterm{def:g6:plants-are-producers:matter}{matéria
orgânica} que os pais precisam obter e reconstruir antes
(\cref{def:g6:matter-from-food:production}). Pagamentos: o \omterm{def:g2:animals-grow-up:born}{ovo} diário
da galinha, a safra de \omterm{prop:g3:flowers-fruits-seeds:transform}{sementes} do ano de fartura da faia, o \omterm{def:g2:animals-grow-up:born}{leite} da
corça (ou os chifres que o cervo faz crescer de novo).
\end{solution}

\begin{solution}{exo:g8:reproduction-and-environment:4}
Todos os indivíduos de uma \omterm{def:g5:biodiversity:species}{espécie} num mesmo lugar. A reprodução (e as
chegadas) acrescentam; as mortes e as partidas subtraem; o alimento, o
espaço, os predadores e as \omterm{def:g6:exploring-our-environment:conditions}{condições} do \omterm{def:g6:exploring-our-environment:components}{ambiente} definem o teto.
\end{solution}

\begin{solution}{exo:g8:reproduction-and-environment:5}
A predação (garças, larvas de libélula); o alimento e o espaço (só
existe tanto para comer e tantos lugares); as \omterm{def:g6:exploring-our-environment:conditions}{condições} (geada, seca,
a lagoa que encolhe).
\end{solution}

\begin{solution}{exo:g8:reproduction-and-environment:6}
A qualidade do \omterm{def:g2:where-animals-live:habitat}{hábitat}, escrita em \omterm{def:g2:animals-grow-up:born}{ovos}: um \omterm{def:g2:where-animals-live:habitat}{hábitat} rico em \omterm{ex:g2:animals-grow-up:butterfly}{lagartas}
constrói uma \omterm{def:g4:animal-reproduction:sexual}{fêmea} capaz de prover dez; um \omterm{def:g2:where-animals-live:habitat}{hábitat} pobre, seis ou
sete.
\end{solution}

\begin{solution}{exo:g8:reproduction-and-environment:7}
A maquinaria reprodutiva das ovelhas é acionada pelos dias que
encurtam; a iluminação artificial repete esse sinal no calendário de
quem cria, e o rebanho pare cedo. Isso prova que o tempo é lido a
partir da própria deixa --- mude a \omterm{def:g6:exploring-our-environment:conditions}{luz} e o ``calendário'' obedece.
\end{solution}

\begin{solution}{exo:g8:reproduction-and-environment:8}
Tempo: a deixa e a desova foram normais --- sem descompasso. Orçamento:
os pais proveram normalmente --- milhares de \omterm{def:g2:animals-grow-up:born}{ovos}. Freios: os três
apertaram --- a seca reduziu o espaço e o \omterm{def:g7:what-is-respiration:respiration}{oxigênio}, o amontoamento
matou de fome, as garças concentraram a caçada. \omterm{def:g2:where-animals-live:habitat}{Hábitat}: a própria
lagoa de reprodução quase sumiu --- a linha decisiva.
\end{solution}

\begin{solution}{exo:g8:reproduction-and-environment:9}
Os filhotes nascem para um banquete que já acabou: as \omterm{ex:g2:animals-grow-up:butterfly}{lagartas}
chegaram ao pico enquanto a deixa de duração do dia ainda dizia
``espere''. O sucesso de ninho cai, embora o tempo, o orçamento e os
predadores estejam bondosos --- a linha do tempo da auditoria falha: a
deixa e o pico do recurso já não coincidem.
\end{solution}

\begin{solution}{exo:g8:reproduction-and-environment:10}
Os milhões são a estratégia: cada \omterm{def:g2:animals-grow-up:born}{ovo} é uma aposta jogada contra a
predação, a fome e as \omterm{def:g6:exploring-our-environment:conditions}{condições}, e os freios recolhem quase todas ---
esse é o plano, e não o fracasso dele. Em média, os sobreviventes
substituem os pais; o gasto \emph{é} como uma estratégia sem cuidado
compra os poucos vencedores dela (o
\cref{exo:g4:animal-reproduction:7} já tinha dito isso).
\end{solution}

\begin{solution}{exo:g8:reproduction-and-environment:11}
Drenar um brejo: remove \omterm{def:g2:where-animals-live:habitat}{hábitat} de reprodução --- a linha do \omterm{def:g2:where-animals-live:habitat}{hábitat}.
Caixas de ninho: acrescentam espaço --- aliviando esse freio.
Pulverizar em maio: esvazia o pico de alimento --- o freio do alimento,
no pior momento. Restaurar um campo alagado: devolve \omterm{def:g2:where-animals-live:habitat}{hábitat} e
alimento juntos (\cref{ex:g5:biodiversity:storks}).
\end{solution}

\begin{solution}{exo:g8:reproduction-and-environment:12}
O sapo fica na ponta dos muitos e sem cuidado da linha das
estratégias: milhares de \omterm{def:g2:animals-grow-up:born}{ovos} baratos, sem guarda nenhuma. A
aritmética da pirâmide garante que quase toda a matéria de girino
precisa alimentar os \omterm{def:g5:food-webs:roles}{consumidores} da lagoa em vez de virar sapo ---
cada andar repassa só uma fração. E os freios são onde essa aritmética
é aplicada: a predação, o alimento, o espaço e as \omterm{def:g6:exploring-our-environment:conditions}{condições} recolhem o
excedente a cada estação, deixando em média um casal para substituir
um casal. A lagoa não fica sólida de sapo porque a estratégia já
planeja as perdas, a pirâmide as precifica e os freios as recolhem.
\end{solution}

\begin{solution}{pb:g8:reproduction-and-environment:1}
\textbf{1.} Os dias que se alongam na primavera dão a deixa da
postura; o tempo escolhido mira os filhotes famintos no pico de
\omterm{ex:g2:animals-grow-up:butterfly}{lagartas} dos carvalhos.

\textbf{2.} A \cref{prop:g8:reproduction-and-environment:cost}: a
reprodução é provida a partir do \omterm{def:g2:where-animals-live:habitat}{hábitat} --- pais construídos por uma
estação rica provêm ninhadas maiores depois dela.

\textbf{3.} Predação sobre \omterm{def:g2:animals-grow-up:born}{ovos} e filhotes; falta de alimento em ondas
de frio; \omterm{def:g6:exploring-our-environment:conditions}{condições} --- uma semana chuvosa e gelada no ninho; (o espaço
é justamente o que as próprias caixas fornecem).

\textbf{4.} Porque os \omterm{def:g2:animals-grow-up:born}{ovos} sozinhos não explicam nada: a auditoria
precisa do encaixe entre deixa e pico (o registro do tempo), do
orçamento (o que a estação ofereceu), dos freios e do \omterm{def:g2:where-animals-live:habitat}{hábitat} (os
serviços do terreno). O balanço só é legível ao lado das \omterm{def:g3:bones-and-muscles:skeleton}{colunas} de
contexto dele.

\textbf{5.} O sucesso de ninho cai: os filhotes nascem depois de o
pico antecipado ter passado --- a linha do tempo: deixa inalterada,
recurso deslocado, coincidência quebrada.

\textbf{6.} Orçamento e freios, por meio do \omterm{def:g2:where-animals-live:habitat}{hábitat}: os \omterm{def:g1:plants-around-us:tree}{galhos} mortos
guardavam os \omterm{prop:g3:sorting-living-things:groups}{insetos} e os buracos naturais, e a faixa ceifada todo mês
não alimentava \omterm{ex:g2:animals-grow-up:butterfly}{lagarta} nem aranha nenhuma --- freio do alimento
apertado e \omterm{def:g2:where-animals-live:habitat}{hábitat} empobrecido, com tempo bondoso ou sem ele.

\textbf{7.} A predação. A ninhada foi provida antes de o pedágio do
gavião cair --- o orçamento da \omterm{def:g4:animal-reproduction:sexual}{fêmea} a definiu; o gavião recolhe
depois, sobre os filhotes. Linhas diferentes do balanço, cobradores
diferentes.

\textbf{8.} O freio do espaço de ninho --- mais lugares de reprodução,
então mais casais no total. Intocado: o teto de alimento por
território, e é por isso que o sucesso de cada caixa se manteve em vez
de subir.

\textbf{9.} ``Dez anos de ninhadas acompanharam o terreno: anos ricos
em \omterm{ex:g2:animals-grow-up:butterfly}{lagarta} e carvalhos menos arrumados encheram as caixas com oito a
dez \omterm{def:g2:animals-grow-up:born}{ovos}, e anos pobres e arrumados demais, com seis --- o tamanho da
ninhada é a qualidade do \omterm{def:g2:where-animals-live:habitat}{hábitat}, escrita em \omterm{def:g2:animals-grow-up:born}{ovos}.''

\textbf{10.} As caixas substituem só o espaço de ninho. Os carvalhos
são alvo da deixa e orçamento ao mesmo tempo --- o pico de \omterm{ex:g2:animals-grow-up:butterfly}{lagartas} a
que o tempo da postura mira e o alimento que provê a ninhada e os
filhotes --- e ainda aliviam um freio (a despensa natural dos \omterm{def:g1:plants-around-us:tree}{galhos}
deles). Cinquenta caixas num estacionamento são prateleiras numa
despensa esvaziada.

\textbf{11.} Deixar a faixa de campo sem ceifa durante a primavera ---
reabastecendo o orçamento de alimento; manter (ou repor) \omterm{def:g1:plants-around-us:tree}{galhos} mortos
e carvalhos velhos --- preservando o \omterm{def:g2:where-animals-live:habitat}{hábitat} e o pico a que o
calendário inteiro mira.

\textbf{12.} De um \omterm{def:g2:where-animals-live:habitat}{hábitat}: a deixa na hora certa, o alimento para
pagar por ela e freios brandos o bastante para deixar, em média, um
casal substituir um casal.
\end{solution}
